% =====================================================================
\chapter{Representaciones 3D y su diferenciabilidad}
\label{tema:1}
% =====================================================================

Antes de optimizar una forma hay que decidir cómo se describe con números. La elección no es neutral: determina qué cambios de forma son posibles (por ejemplo, si pueden aparecer agujeros), cuánta memoria hace falta y, sobre todo, si la magnitud que se quiere optimizar es una función derivable de los parámetros. Este tema presenta el esquema común de todos los optimizadores del módulo y repasa las representaciones clásicas y neuronales siguiendo la clase «Classical 3D Geometry Representations» del curso CS348n de Stanford~\cite{cs348n}, el artículo de DeepSDF~\cite{deepsdf}, el de \emph{marching cubes}~\cite{mc} y el de las codificaciones con tablas \emph{hash} de Müller \emph{et al.}~\cite{ingp}.

\section{El optimizador geométrico}

\begin{definicion}[Optimizador geométrico]\label{def:opt-geom}
Sea $\Theta\subseteq\R^n$ un conjunto de parámetros de diseño. Un \emph{optimizador geométrico} resuelve
\[
\min_{\theta\in\Theta}\ J(\theta)=\mathcal{L}\bigl(S(G(\theta))\bigr)\quad\text{sujeto a}\quad g_j(\theta)\le 0,\ j=1,\dots,m,
\]
donde $G$ construye una geometría a partir de $\theta$, $S$ la evalúa (un simulador físico, que devuelve campos como desplazamientos o temperaturas, o un renderizador, que devuelve imágenes), $\mathcal{L}$ resume la evaluación en un número (flexibilidad, calor disipado, error frente a fotografías) y las $g_j$ expresan restricciones de fabricación, de volumen o de presupuesto.
\end{definicion}

\addtocounter{figure}{1}%

El gradiente $\nabla_\theta J$ puede obtenerse de tres maneras. Las diferencias finitas (módulo~4, tema~6) necesitan $n+1$ evaluaciones de $S$ y son inviables cuando $n$ es del orden de miles. El modo inverso de la diferenciación automática o, cuando $S$ resuelve un sistema de ecuaciones, el método adjunto (módulo~10, §5.4, y \cref{tema:2}) calculan el gradiente completo con un coste comparable a una o dos evaluaciones, independientemente de $n$. Y cuando $S$ no es derivable o es demasiado caro, se aprende un sustituto $\hat s$ con unas decenas de evaluaciones y se optimiza sobre él (\cref{tema:3}). En los tres casos, la primera condición es que la geometría $G(\theta)$ dependa de $\theta$ de forma suave, y eso depende de la representación.

\section{Representaciones explícitas: mallas, nubes de puntos y vóxeles}

\begin{definicion}[Malla triangular]\label{def:malla}
Una \emph{malla triangular} es un par $(V,F)$ formado por una lista de vértices $v_1,\dots,v_{\abs{V}}\in\R^3$ y una lista de caras, cada una un triplete ordenado $(a,b,c)$ de índices de vértices. El orden de los índices define la orientación de la cara y su normal $n_f=(v_b-v_a)\times(v_c-v_a)/\norm{(v_b-v_a)\times(v_c-v_a)}$. El conjunto $E$ de aristas está formado por los pares no ordenados de vértices consecutivos de alguna cara. La malla es \emph{cerrada} si cada arista pertenece exactamente a dos caras, y está \emph{orientada de forma coherente} si cada arista compartida se recorre en sentidos opuestos en sus dos caras.
\end{definicion}

\begin{proposicion}\label{prop:3f2e}
En una malla triangular cerrada, $2\abs{E}=3\abs{F}$. En particular, el número de caras es par.
\end{proposicion}

La característica de Euler $\chi=\abs{V}-\abs{E}+\abs{F}$ de una malla cerrada y conexa que no se corta a sí misma depende solo de su topología: $\chi=2-2g$, donde $g$ es el género o número de «asas» (0 para una esfera, 1 para un toro). Este resultado de topología se admite sin demostración.

\begin{corolario}\label{cor:euler}
Si una malla triangular cerrada tiene $\chi=2$ (topología de esfera), entonces $\abs{F}=2\abs{V}-4$ y $\abs{E}=3\abs{V}-6$.
\end{corolario}

El corolario da una regla de memoria útil: una malla cerrada tiene aproximadamente el doble de caras que de vértices. Con posiciones en coma flotante de 32 bits (12 bytes por vértice) e índices enteros de 32 bits (12 bytes por cara), cuesta unos 36 bytes por vértice.

\begin{definicion}[Nube de puntos y distancia de Chamfer]\label{def:chamfer}
Una \emph{nube de puntos} es un conjunto finito $A\subset\R^3$, opcionalmente con normales. La \emph{distancia de Chamfer} entre dos nubes $A$ y $B$ es
\[
\CD(A,B)=\frac{1}{\abs{A}}\sum_{a\in A}\min_{b\in B}\norm{a-b}^2+\frac{1}{\abs{B}}\sum_{b\in B}\min_{a\in A}\norm{a-b}^2 .
\]
\end{definicion}

La distancia de Chamfer es derivable respecto a las coordenadas de los puntos salvo en los empates del mínimo: el gradiente respecto a $a$ empuja $a$ hacia su vecino más próximo en $B$. Es la pérdida habitual para ajustar nubes de puntos~\cite{cmu16825}. Una \emph{rejilla de vóxeles} de resolución $n$ guarda un valor (ocupación, densidad o distancia) en cada una de las $n^3$ celdas de una caja; es la representación más sencilla de derivar, porque cada parámetro es un valor de celda, pero su memoria crece con el cubo de la resolución.

\begin{proposicion}[Volumen de una malla cerrada y su gradiente]\label{prop:volumen}
Sea $(V,F)$ una malla triangular cerrada, orientada de forma coherente con normales hacia fuera, que encierra la región $\Omega$. Entonces
\[
\Vol(\Omega)=\frac16\sum_{(a,b,c)\in F}\det(v_a,v_b,v_c)=\frac16\sum_{(a,b,c)\in F}v_a\cdot(v_b\times v_c),
\]
y la derivada respecto a la posición de un vértice $v_i$ es $\nabla_{v_i}\Vol=\frac16\sum v_b\times v_c$, sumando sobre las caras que contienen a $i$, escritas en el orden cíclico $(i,b,c)$.
\end{proposicion}

\addtocounter{definicion}{1}%

\section{Representaciones implícitas y funciones de distancia con signo}

\begin{definicion}[Superficie implícita y SDF]\label{def:sdf}
Una función continua $f:\R^3\to\R$ define el sólido $\Omega=\{f\le 0\}$ y la superficie implícita $\{f=0\}$. Se dice que $f$ es una \emph{función de distancia con signo} (SDF) de $\Omega$ si $f(x)=-\operatorname{dist}(x,\partial\Omega)$ para $x\in\Omega$ y $f(x)=\operatorname{dist}(x,\partial\Omega)$ fuera de $\Omega$.
\end{definicion}

Por ejemplo, $f(x)=\norm{x-c}-r$ es la SDF de la bola de centro $c$ y radio $r$, y en todo punto $x\ne c$ su gradiente $\nabla f=(x-c)/\norm{x-c}$ tiene norma~1.

\begin{proposicion}[Una SDF es 1-lipschitziana]\label{prop:lipschitz}
Si $f$ es la SDF de $\Omega$, entonces $\abs{f(x)-f(y)}\le\norm{x-y}$ para todo $x,y$. En consecuencia, en los puntos en que $f$ es derivable, $\norm{\nabla f}\le 1$.
\end{proposicion}

De hecho, una SDF cumple casi en todas partes la \emph{ecuación eikonal} $\norm{\nabla f}=1$, y su gradiente sobre la superficie es la normal exterior. Las operaciones booleanas son inmediatas: si $f_1$ y $f_2$ describen $\Omega_1$ y $\Omega_2$, entonces $\min(f_1,f_2)$ describe la unión, $\max(f_1,f_2)$ la intersección y $\max(f_1,-f_2)$ la diferencia $\Omega_1\setminus\Omega_2$. Los conjuntos son exactos, aunque los valores resultantes pueden ser solo cotas de la distancia (por ejemplo, dentro de una unión). Esta robustez de las operaciones es el argumento del modelado implícito que usan herramientas comerciales como nTop, que publica que sus operaciones booleanas, desfases y redondeos no fallan, y que modula celosías y espesores con campos escalares procedentes de simulaciones~\cite{ntop}.

Para fabricar o renderizar con rasterización hay que convertir $f$ en una malla. El algoritmo de \emph{marching cubes}~\cite{mc} evalúa $f$ en los nodos de una rejilla y, en cada arista con cambio de signo entre $f_0<0$ y $f_1>0$, coloca un vértice por interpolación lineal en la fracción $t=f_0/(f_0-f_1)$ de la arista; una tabla indexada por el patrón de signos de los ocho vértices del cubo da los triángulos. El coste es una evaluación de $f$ por nodo, $n^3$ en total, y el error de posición es del orden del paso de la rejilla.

\section{Campos neuronales}

Un \emph{campo neuronal} representa $f$ con una red, $f_\varphi(x)$, típicamente un perceptrón multicapa (módulo~9). DeepSDF~\cite{deepsdf} ajusta una red de ocho capas de 512 unidades a valores de SDF muestreados con la pérdida $L^1$ recortada
\[
\ell\bigl(f_\varphi(x),s\bigr)=\bigl\lvert\clip(f_\varphi(x),-\delta,\delta)-\clip(s,-\delta,\delta)\bigr\rvert,
\]
con $\delta=0{,}1$, que concentra la capacidad cerca de la superficie. Para representar una familia de formas, DeepSDF asigna a cada forma un código latente $z_i$ que se optimiza junto con los pesos, con un término $\norm{z_i}^2/\sigma^2$ que procede de un prior gaussiano (un \emph{autodecodificador}); a una forma nueva se le asigna su código optimizando solo $z$ con los pesos congelados. La analogía con un LLM es directa: los códigos $z_i$ desempeñan el papel de las filas de la matriz de embeddings (módulo~12, §4.1), que también son parámetros libres aprendidos por descenso de gradiente.

Cuando no se dispone de valores de SDF sino solo de puntos de la superficie, se impone $f_\varphi=0$ en ellos y se añade la regularización eikonal $\E_x\bigl[(\norm{\nabla_xf_\varphi(x)}-1)^2\bigr]$, cuyo gradiente respecto a $x$ se obtiene por diferenciación automática~\cite{cs348nhw2}. Los perceptrones con entradas $(x,y,z)$ aprenden mal los detalles finos; dos remedios son la codificación posicional sinusoidal (\cref{tema:4}), análoga a la del Transformer (módulo~14, §4.2), y las rejillas multirresolución de vectores entrenables indexadas con funciones \emph{hash}, con las que Müller \emph{et al.} entrenan campos neuronales de calidad en segundos~\cite{ingp}.

\section{Diferenciabilidad y elección de la representación}

\begin{table}[h]
\centering\small
\begin{tabularx}{\textwidth}{@{}l>{\raggedright\arraybackslash}X>{\raggedright\arraybackslash}Xr@{}}
\toprule
Representación & Parámetros derivables & Cambios de topología & Memoria \\
\midrule
Malla & posiciones de los vértices (conectividad fija) & no, salvo remallado & 4,1 MB \\
Nube de puntos & coordenadas de los puntos & no aplica (sin superficie) & --- \\
Vóxeles (ocupación) & valor de cada celda & sí & 0,29--2,3 MB \\
SDF en rejilla & valor de cada nodo & sí & 9,2 MB \\
Campo neuronal & pesos de la red & sí & 7,4 MB \\
\bottomrule
\end{tabularx}
\caption{Comparación para un soporte que cabe en una caja de $120\times 60\times 40$ mm con resolución de 0,5 mm. Vóxeles: $240\cdot 120\cdot 80=2\,304\,000$ celdas (1 bit o 1 byte por celda). SDF: 4 bytes por nodo. Malla: superficie de la caja $28\,800$ mm² en triángulos rectángulos de catetos 0,5 mm, $230\,400$ caras y $115\,200$ vértices. Red: 8 capas de 512 unidades, $1{,}84\cdot 10^{6}$ parámetros de 4 bytes.}
\label{tab:repr}
\end{table}

El \cref{tab:repr} resume el compromiso. Una malla con conectividad fija es compacta y es lo que consumen los motores y las impresoras, pero su optimización no puede abrir agujeros; las representaciones volumétricas (densidades, SDF, campos) sí pueden, y por eso la optimización de topología (\cref{tema:2}) trabaja con densidades por elemento. En todas ellas la derivabilidad se refiere a la magnitud que se evalúa: el volumen de una malla es un polinomio en las coordenadas (\cref{prop:volumen}), pero la imagen rasterizada de esa misma malla es constante a trozos (\cref{tema:4}).

\section{Caso de estudio: ajuste de una esfera a un escaneo parcial}

\addtocounter{definicion}{1}%

\begin{table}[h]
\centering\small
\begin{tabularx}{\textwidth}{@{}l>{\raggedright\arraybackslash}X>{\raggedright\arraybackslash}p{2.6cm}>{\raggedright\arraybackslash}X@{}}
\toprule
Herramienta & Para qué sirve & Licencia & Cuándo elegirla \\
\midrule
trimesh & cargar, reparar y medir mallas (volumen, estanqueidad) & MIT & validación y utilidades en Python puro \\
Open3D & nubes de puntos, registro, reconstrucción & MIT & datos de escáner y visualización \\
PyMeshLab & filtros de MeshLab (remallado, simplificación) & GPL-3.0 & procesado por lotes con algoritmos maduros \\
libigl & geometría discreta (laplacianos, parametrización) & MPL-2.0 & investigación y prototipos de procesado de geometría \\
scikit-image & \emph{marching cubes} sobre rejillas & BSD-3 & extraer mallas de SDF o densidades \\
PyTorch3D & mallas y nubes por lotes, pérdidas y renderizado derivables & BSD & ajuste de formas con gradientes en PyTorch \\
Kaolin (NVIDIA) & operadores 3D derivables y conversiones & Apache-2.0 & ecosistema NVIDIA con GPU \\
\bottomrule
\end{tabularx}
\caption{Tecnologías del tema~1. Licencias según los metadatos publicados por cada proyecto.}
\label{tab:tec1}
\end{table}

\addtocounter{definicion}{1}%

